% ============================================================
% git-instructions-main.tex — A "GROUPING" SUBDOCUMENT
% This file doesn't contain content of its own — it just
% bundles the individual Git-related subfiles together in
% order, so main.tex only needs one \subfile call to pull in
% the whole Git section.
%
% IMPORTANT: even though this file sits "between" main.tex and
% the individual topic files, its \documentclass line still
% points to main.tex — NOT to itself. Every subfile in a nested
% chain points to the same ultimate root document.
% ============================================================

\documentclass[../main.tex]{subfiles}

\begin{document}

\section{Git and GitHub}
% This is now the ONLY \section command in the whole Git
% bundle. Everything pulled in below must use \subsection
% (not \section) so it nests under this one TOC entry
% instead of creating its own.

\begingroup
  % \begingroup ... \endgroup creates a "scope" — any command
  % redefinitions inside only last until \endgroup, so this
  % trick doesn't affect anything outside this block.
  
	\let\originalsubsection\subsection
  	\let\originalsubsubsection\subsubsection
  	% Save the REAL commands under safe temporary names first,
  	% before we overwrite \section and \subsection below.
  

	\renewcommand{\section}[1]{\originalsubsection{#1}}
  	\renewcommand{\subsection}[1]{\originalsubsubsection{#1}}
  	% consistent.

	\subfile{download-git-bash.tex}
	\subfile{github-ssh.tex}
	\subfile{git-init-existing-folder.tex}

\endgroup

\end{document}